% =====================================================================
%  A characteristic outflow posed on the relaxed state
%  Discussion note for a conversation on NSCBC beyond the frozen gas:
%  flashing / relaxing two-phase flow at a subsonic outlet, with flow
%  reversal.  Reuses docs/PLAN_relaxed_boundary.md (derivation, decisions,
%  measurements) and camr_ps_model.tex S8 (the ghost construction).
%  Build: pdflatex nscbc_relaxed_discussion.tex (twice).
% =====================================================================
\documentclass[11pt,a4paper]{article}
\usepackage[margin=25mm]{geometry}
\usepackage{amsmath,amssymb}
\usepackage{booktabs}
\usepackage{array}
\usepackage[colorlinks=true,linkcolor=blue,citecolor=blue,urlcolor=blue]{hyperref}
\usepackage{xcolor}
\newcommand{\dd}[2]{\frac{\partial #1}{\partial #2}}
\newcommand{\code}[1]{\texttt{\small #1}}
\newcommand{\Lin}{\mathcal{L}_{\mathrm{in}}}
\newcommand{\W}{\mathcal{W}}
\newcommand{\uout}{u_{\mathrm{out}}}
\newcommand{\Pinf}{P_\infty}
\newcommand{\ceq}{c_{\mathrm{eq}}}
\newcommand{\cf}{c_{f}}
\newcommand{\cw}{c_{w}}
\newcommand{\uvec}{\mathbf{u}}
\newcommand{\Ivec}{\mathsf{I}}

\title{A characteristic outflow posed on the relaxed state\\
\large Notes for a discussion of NSCBC in flashing, relaxing two-phase flow}
\author{M.~Day \quad (CAMR, branch \code{co2-eos})}
\date{September 2026 --- work in progress (tree \code{7abcc31})}

\begin{document}
\maketitle

\begin{abstract}
\noindent
We are simulating CO$_2$ pipeline depressurisation with a six-equation
two-phase model (Pelanti--Shyue) on real-fluid equations of state
(Peng--Robinson and GERG-2008), whose closures are an instantaneous
pressure-equilibrium constraint between the phases, finite-rate thermal
and mass-transfer relaxation, and nucleation (flash).  The
interior scheme is wave propagation.  At a subsonic outflow the
classical NSCBC construction, written on the frozen (Wallis)
eigenstructure, reflects a large fraction of any wave that has been
relaxed: the relaxed wave has components on both frozen
characteristics, and any condition that sets the entering frozen wave
from the far field returns part of the leaving relaxed one.  We built an
alternative that treats the strength of the entering discrete wave as
the one modelled quantity and poses the boundary condition on the state
the cell relaxes to at the end of the step, with the relaxed impedance
measured from the trial evaluations themselves.  On plane relaxed waves
it takes the reflection from $0.24$ to $0.01$ with nothing lost on pure
vapour.  On the production outlet --- a supersonic two-phase jet core
impinging on the outflow plane with the surrounding face in reversed
flow --- it has taught us, run by run, which parts of the classical
construction were doing what.  This note frames the construction against
Poinsot--Lele for a reader whose reference application is reacting flow,
and lists the questions we would like to discuss.
\end{abstract}

% ---------------------------------------------------------------------
\section{Where the applications overlap}
% ---------------------------------------------------------------------

The reacting-flow and the flashing-flow outlet share the structure that
made NSCBC necessary in the first place.  Both have a subsonic outflow
through which acoustic waves must leave without reflection while the
mean pressure is held to a far-field value; both carry a convected
scalar field (composition; volume fraction and phase state) that must
cross the boundary without generating pressure; both have a stiff source
term whose time scale is far below the acoustic one.  Sutherland and
Kennedy's observation for reacting flows --- that the source belongs
\emph{in} the boundary relation, because the characteristic amplitudes
are not decoupled from a stiff source over a time step --- is the
starting point here too.

Both applications also meet the two situations the literature treats
least well.  Reversed flow on part of the outlet (a recirculating jet, a
Mach-disk recompression) turns an outflow face into a partial inflow
with the entering entropy and vorticity waves undetermined.  And the
finite-volume setting has no boundary node: the condition must be
imposed through a ghost state or on a face flux, so the drift and the
near-sonic ill-conditioning that Masset et al.\ analyse for ghost-cell
CBC are shared.


% ---------------------------------------------------------------------
\section{The model, in the variables used below}
% ---------------------------------------------------------------------

The six-equation model of Pelanti and Shyue carries, for two phases
$k = 1$ (liquid) and $k = 2$ (vapour, or supercritical fluid) sharing
one velocity $\uvec$, the volume fraction $\alpha_1$
($\alpha_2 = 1 - \alpha_1$), the partial masses $m_k = \alpha_k\rho_k$,
the mixture momentum $\rho\uvec$ and the phase total energies
$\mathcal{E}_k = \alpha_k\rho_k E_k$ with $E_k = e_k + \tfrac12|\uvec|^2$.
Each phase has its own density $\rho_k$, internal energy $e_k$,
temperature $T_k$ and pressure $P_k = P_k(\rho_k, e_k)$ from its own
branch of the equation of state; the mixture pressure is
$P = \alpha_1 P_1 + \alpha_2 P_2$ and the mixture density
$\rho = m_1 + m_2$.  The system is
\begin{align}
  \partial_t\alpha_1 + \uvec\!\cdot\!\nabla\alpha_1 &= \mathcal{S}_\alpha, \\
  \partial_t m_k + \nabla\!\cdot\!(m_k\uvec) &= \pm\Gamma, \\
  \partial_t(\rho\uvec) + \nabla\!\cdot\!(\rho\uvec\otimes\uvec + P\Ivec) &= 0, \\
  \partial_t\mathcal{E}_k + \nabla\!\cdot\!(\mathcal{E}_k\uvec + \alpha_k P_k\uvec)
    &= \mp P_I\,\partial_t\alpha_1 \pm \mathcal{Q} \pm \Gamma H_I ,
\end{align}
with $\mathcal{Q}$ the thermal relaxation ($T_1 \to T_2$), $\Gamma$ the
mass transfer (driving the Gibbs energies together; nucleation of a
phase that is absent is a separate flash operator), the interfacial
pressure $P_I = \tfrac12(P_1 + P_2)$, and $\mathcal{S}_\alpha$ the term
that couples the two phase pressures.  In Pelanti and Shyue's
formulation $\mathcal{S}_\alpha = \mu\,(P_1 - P_2)$ with $\mu \to
\infty$; we take that limit literally rather than as a rate.  The
hyperbolic step advances the six equations with $\alpha_1$ advected and
each phase carrying its own pressure, and after every step $\alpha_1$
is reset, cell by cell, to the value at which $P_1 = P_2$ (a Newton on
$\alpha_1$, with the interfacial work $\mp P_I\,\delta\alpha_1$
deposited in the phase energies as the discrete form of the
non-conservative term).  Pressure equality is therefore a constraint of
the closure --- mechanical equilibrium, imposed by redistributing
volume between the phases --- and not a relaxation process with a time
scale: a finite-rate version was measured inert and removed.  Thermal
and mass-transfer relaxation are finite-rate in the production closure
(the coupled ``X3'' operator, which integrates both on the $P_1 = P_2$
manifold and whose fixed point is the equilibrium state).  The mixture
total energy $\rho E = \mathcal{E}_1 + \mathcal{E}_2$ is fluxed in
strict conservation form; the volume-fraction and phase-energy
equations are not in divergence form, which is why the interior uses
wave propagation rather than flux differencing.

The phases are evaluated on real-fluid equations of state: Peng--Robinson
(analytic, and a tabulated form) and GERG-2008 (the reference
multiparameter equation, tabulated), each with a branch-locked
liquid/vapour inversion so that a phase is always queried on its own
branch and ``no root'' is reported rather than substituted.  Every
saturation quantity the closures use ($P_{\mathrm{sat}}(T)$,
$T_{\mathrm{sat}}(P)$, the saturated phase states) is the equation of
state's own coexistence curve.

\paragraph{The homogeneous-equilibrium limit.}  As all three relaxation
times go to zero, $P_1 = P_2$, $T_1 = T_2$ and $g_1 = g_2$ at every
instant; $\alpha_1$, $m_1$ and $m_2$ cease to be independent and are
fixed by $(\rho, e)$ through the saturation curve, and the six equations
collapse to the three-equation Euler system for the mixture with an
equilibrium (HEM) equation of state $P = P_{\mathrm{HEM}}(\rho, e)$.  The
HEM limit is what the closures relax toward and what our exact
two-phase reference solutions are computed on; the finite-rate model
sits between the frozen system and that limit.

\paragraph{Frozen and relaxed sound speeds.}  Throughout this note
``frozen'' means that the relaxation terms are ignored: the sources
$\mathcal{S}_\alpha$, $\mathcal{Q}$ and $\Gamma$ are set to zero, so each
phase keeps its own pressure, temperature and mass, and the mixture is
two phases compressed side by side.  Three sound speeds follow from
which of the closures act on a disturbance.  With none of them acting
--- the frozen system --- each phase is compressed along its own
isentrope at fixed volume fraction, the mixture's bulk modulus is the
mass-weighted sum of the phase moduli, and the speed is
$\cf^2 = Y_1 c_1^2 + Y_2 c_2^2$ with $Y_k = m_k/\rho$ (Wallis).  In the
model as we solve it no physical wave travels at $\cf$: the pressure
constraint is re-imposed after every step, so $\cf$ is the speed of the
characteristics of the hyperbolic operator within a step, the object
the split scheme propagates and the Riemann solver is built on, and
nothing more.  With the pressure constraint acting, the volume
fractions redistribute so that the more compressible phase takes more
of the volume change; the mixture's compressibility is the
volume-weighted sum of the phase compressibilities,
$1/(\rho\cw^2) = \sum_k \alpha_k/(\rho_k c_k^2)$, and the speed $\cw$
(Wood) satisfies $\cw \le \cf$, with the gap large in a bubbly mixture
(here $\cw = 174$~m/s against $\cf = 282$~m/s); this is the speed of a
wave in the model with the thermal and mass-transfer legs switched off.
With those legs also in equilibrium, evaporation and condensation along
the wave absorb part of the pressure change and the speed is lower
still: $\ceq$, the sound speed of the homogeneous equilibrium
mixture.  Two further symbols are used below.  The
\emph{impedance} $Z = \rho c$ is the ratio $\Delta P/\Delta u$ across an
acoustic wave travelling at speed $c$, so $Z_f = \rho\cf$ for a frozen
wave and $Z_{\mathrm{rel}} = \rho\,c_{\mathrm{rel}}$ for a relaxed one;
and $c_{\mathrm{rel}}$ denotes the speed at which a wave actually travels
in the relaxed flow --- $\cw$ when only the pressure constraint has
acted, $\ceq$ when the thermal and mass-transfer legs are in
equilibrium too, and in between for the finite-rate closure (139~m/s
measured for the production closure on the mixture above).  The scheme
propagates frozen characteristics within the step and applies the
constraint and the relaxation legs afterwards; the physical waves
travel at $c_{\mathrm{rel}}$.

% ---------------------------------------------------------------------
\section{Where flashing flow differs from reacting flow}
% ---------------------------------------------------------------------

Three differences matter at the boundary, and they compound.

\paragraph{The source changes the eigenstructure.}  A reaction term
moves the state along the composition direction and leaves the acoustic
eigenstructure of the gas essentially alone.  Relaxation in a two-phase
mixture does not: the frozen (Wallis) sound speed of the six-equation
system,
\begin{equation}
  \cf^2 = Y_1 c_1^2 + Y_2 c_2^2 ,
\end{equation}
is the speed of the characteristics the scheme propagates, while the
speed at which a wave actually travels once the pressure constraint has
acted is the Wood speed,
\begin{equation}
  \frac{1}{\rho\,\cw^2} = \frac{\alpha_1}{\rho_1 c_1^2} + \frac{\alpha_2}{\rho_2 c_2^2},
\end{equation}
and after thermal and mass-transfer relaxation the equilibrium speed
$\ceq$, lower still.  The sub-characteristic condition
$\ceq \le \cw \le \cf$ is what makes the split scheme stable, and it is
also what makes the boundary reflect.  A relaxed wave travelling at
$c_{\mathrm{rel}}$ carries $\Delta P/\Delta u = Z_{\mathrm{rel}} =
\rho\,c_{\mathrm{rel}}$; the frozen Riemann invariants of the boundary
cell, $R^\pm = u \pm P/Z_f$, are built with $Z_f = \rho\cf$, so that
same wave has a non-zero component on the \emph{entering} frozen
characteristic,
\begin{equation}
  \Delta R^- = \Delta P\Big(\frac{1}{Z_{\mathrm{rel}}} - \frac{1}{Z_f}\Big),
\end{equation}
which any condition written on the frozen characteristics sends back
into the domain.  The reflected fraction of the pressure amplitude is
the impedance mismatch $(Z_f - Z_{\mathrm{rel}})/(Z_f + Z_{\mathrm{rel}})$.
For our saturated mixture at $\alpha_1 = 0.05$, $\cf = 282$~m/s and the
measured $c_{\mathrm{rel}} = 139$~m/s give a predicted reflection of
about $0.3$; measured $0.22$--$0.24$, the same for the NSCBC ghost and
for a plain zero-gradient copy.  Nothing in the invariant construction
is at fault; the basis is.

\paragraph{The source acts along the wave, not only in the cell.}  A
depressurising liquid expands through an isentropic fan along which it
crosses the saturation dome and flashes; the fan's own sound speed
collapses as vapour forms, and the physical exit chokes at the
equilibrium sound speed, far below the frozen one.  Our boundary
construction for a venting face follows exactly that path.  Starting
from the phase states of the boundary cell, it marches the pressure
from the cell value toward the ambient value in substeps (each limited
to a 5\,\% density change of either phase --- a march in pressure along
the isentrope, not a march in time through the CFD step), moves each
phase along its own isentrope at every substep, and, whenever the
mixture enthalpy at the current pressure lies inside the saturation
dome, re-closes the mixture to the equilibrium (HEM) saturated state at
that pressure, which is where the flash enters and where the path's
sound speed falls.  The march stops at the ambient pressure or at the
point of maximal mass flux $\rho u$ along the path, the critical-flow
condition; the state reached is the ghost state.  We measure this
construction on a one-dimensional tube of supercritical liquid CO$_2$
(100~bar, 320~K) venting through one end to 1~bar, scoring the mass
that has left the tube by $80~\mu$s.  The reference is the same tube
computed on a domain twice as long, with the vent plane an interior
interface and 1~bar vapour beyond it, so that no boundary model is
involved; the fan construction vents 4.4\,\% more mass than that
reference.  A relaxation operator that acts on the boundary cell
\emph{after} a frozen wave has entered cannot reproduce this.  The wave
that enters the cell is a frozen HLLC jump; a frozen jump from 100~bar
to about 1~bar on a liquid is an isentropic expansion by 99~bar with no
phase change, and the liquid branch of the equation of state has no
state there, so the trial evaluation refuses it.  The physical path
flashes on the way down, with the pressure falling through the dome in
a two-phase mixture, which is what the substepped march follows and a
single jump cannot.  This is the one place the new construction has had
to concede ground, and it is one of the questions below.

\paragraph{The boundary state has a thermodynamic admissibility
problem.}  Each phase is evaluated on its own branch of a cubic
equation of state, and a boundary value obtained by extrapolation or by
construction may fall outside the region where that phase-locked
evaluation has a solution --- a liquid density and energy for which the
liquid branch has no root, say --- and is then simply not a state.
Masset et al.\ meet the same problem for multicomponent gases with
complex state laws and answer it by shrinking the boundary gradient
until the ghost becomes admissible.  Our implementation avoids floors
and thresholds of that kind: where a construction generates a
non-physical value we do not adjust it, we discard it, count the
refusal, and fall back to a lower-accuracy construction that is
admissible by design.  For example, an extrapolated ghost state whose
phases have no root on their branches is replaced by a copy of the
boundary cell, which is admissible because the boundary cell is; and a
trial state of the relaxed boundary that refuses is replaced by the
wave the interior fan had already sent.

% ---------------------------------------------------------------------
\section{The interior scheme, and the boundary face in its language}
% ---------------------------------------------------------------------

The interior is Berger--LeVeque wave propagation.  At every face the
HLLC fan gives three waves $\W_l$ with speeds $s_l$, and the two cells
that share the face are each updated with the sum of the fluctuations
$\sum s_l \W_l$ of the waves travelling into them, i.e.\ arriving from
upwind; a limited second-order correction flux is added on top.  The
non-conserved quantities --- the volume fraction and the two phase
energies --- are updated by those same fluctuations rather than by a
flux divergence, which would require a carefully synchronised
subtraction of the non-conservative source term.  The viscous
contribution is a Newtonian stress with a constant effective viscosity
$\mu$ (a deck parameter chosen to give the orifice shear layer a
resolved thickness), added in flux form with the viscous heating
partitioned to the phase energies by mass fraction.  After the
hydrodynamic increment, a reaction operator $\mathcal{R}$ acts cell by
cell: the mixture-consistency resyncs, the phase-removal operators (a
phase whose volume fraction or partial mass has fallen below what the
model can resolve is removed and its mass and energy assigned to the
other phase), the pressure-equilibrium projection, the finite-rate
thermal and mass-transfer relaxation, and nucleation.

At a domain face $f$ between the boundary cell $N$ and its ghost $G$
the fan decomposes the jump $q_G - q_N$ that the ghost fill created.  For
subsonic outflow exactly one wave enters cell $N$, the acoustic one; for
subsonic reversed flow two enter, the acoustic wave and the contact,
which carries the ghost's composition and entropy inward.  By the time
the level's fluxes are formed, the boundary cell's predicted increment
is known to the last term --- every interior face, the corrections, the
transverse and viscous contributions, the non-conservative deposit ---
except for the entering acoustic wave:
\begin{equation}
  \Delta q_N = \underbrace{\Delta q^{\mathrm{int}} + \Delta q^{f,\mathrm{out}} + \Delta q^{f,\mathrm{cont}}}_{\text{known}}
  \;-\; \frac{\Delta t}{\Delta x}\, s_{\mathrm{ac}} \W_{\mathrm{ac}} .
  \label{eq:split}
\end{equation}
The construction below replaces $\W_{\mathrm{ac}}$ and leaves everything
else --- including the ghost's contact wave in reversed flow --- as the
scheme produced it.

% ---------------------------------------------------------------------
\section{The construction}
% ---------------------------------------------------------------------

\subsection{The modelled wave}

The entering acoustic wave is taken to be the HLLC star jump on the
$N$ side of $f$ for a prescribed pressure amplitude $a$ [Pa].  Two
speeds of the frozen fan enter here, and they are frozen quantities
because the fan is: $s_{\mathrm{ac}}$ is the HLLC estimate of the speed of
the acoustic wave entering the domain (on a high face
$S_L = \min(u_N - \cf{}_{,N},\, u_G - \cf{}_{,G})$, the outer wave speed on
the domain side), and $c_{\mathrm{eff}} = |s_{\mathrm{ac}} - u_{n,N}|$ is that
wave's speed relative to the boundary cell, which equals $\cf{}_{,N}$
whenever the domain side sets the estimate.  The modelled wave is thus
a frozen wave, as every wave the scheme propagates is; the relaxation
acts on the cell afterwards, in the trial evaluation.  With
$n_{\mathrm{sgn}} = \pm 1$ the outward sign,
\begin{equation}
  S_M(a) = u_{n,N} - \frac{n_{\mathrm{sgn}}\, a}{\rho_N\, c_{\mathrm{eff}}},
  \qquad
  \hat\W(a) = U^\star_N\big(s_{\mathrm{ac}}, S_M(a)\big) - q_N ,
  \label{eq:wave}
\end{equation}
so that the HLLC pressure behind the wave is exactly $P_N + a$ and the
velocity jump is that of an inward-running acoustic wave,
$\mathrm{d}\uout = -a/\rho_N c_{\mathrm{eff}}$.  The star state is built by
the interior's own routine, so the wave satisfies the conservation
identities the scheme relies on, refuses (rather than clamps) when the
jump is inadmissible, and needs no new discretisation at the boundary.
$a = a_{\mathrm{fan}}$ reproduces the wave the ghost fan actually sent; $a
= 0$ is no wave.  The modelled wave enters first order --- formally: on
an outflow the entering wave is $O(r\,p')$ small, and the order series of
\S\ref{sec:meas} finds no trace of it.

\subsection{The trial state}

For a candidate $a$ the boundary cell's end-of-step state and its
relaxed image are
\begin{equation}
  q^{\mathrm{h}}(a) = q_N + \Delta q^{\mathrm{known}} - \frac{\Delta t}{\Delta x}\, s_{\mathrm{ac}}\hat\W(a),
  \qquad
  \tilde q(a) = \mathcal{R}\big(q^{\mathrm{h}}(a);\,\Delta t_r\big),
\end{equation}
where $\mathcal{R}$ is the production reaction chain applied to the
boundary strip only, with its counters snapshotted so the trials leave
no trace in the solution's accounting.  From $\tilde q$ we read the
six-equation mixture pressure $\tilde P = \alpha_1 P_1 + \alpha_2 P_2$,
the outward velocity $\tilde\uout$ and the mixture density $\tilde\rho$.
This is the Sutherland--Kennedy move --- the stiff source inside the
boundary relation --- carried out by evaluating the source on trial
states rather than by deriving a relaxed eigenstructure.

\subsection{The condition}

Let $I(q) = P(q) - Z\,\uout(q)$ be the incoming relaxed invariant of the
boundary cell.  In the LODI frame,
$\partial_t(P - \rho c\,\uout) = -\Lin$ exactly, so every classical
outflow is a statement about $\Lin$, and we discretise it on the relaxed
step:
\begin{align}
  \text{rate form:}\quad & I(\tilde q(a)) - I(q_N) + \Delta t\,K\,(P_N - \Pinf) = 0, \label{eq:rate}\\
  \text{value form:}\quad & I(\tilde q(a)) - \Pinf = 0 . \label{eq:value}
\end{align}
The rate form with $K = 0$ is Poinsot--Lele's $\Lin = 0$; with
$K = \sigma(1-M^2)c/L$ it is the Rudy--Strikwerda anchoring.  The value
form is Hedstrom's: the incoming invariant \emph{equals} its far-field
value every step.  The two differ only when $I$ has drifted from
$\Pinf$; the value form restores it within the step, the rate form
never does.  Both are non-reflecting for an outgoing wave that keeps
$I$ constant --- which a relaxed wave does when $Z$ is its impedance.

\subsection{The impedance}

$Z$ was first taken as the Wood impedance $\rho_N\cw$ of the boundary
cell.  With the wave at the Wood speed (the pressure constraint only,
the thermal and mass-transfer legs off) the condition is then exact to
the single-phase level; with the production
closure the wave travels at the equilibrium speed and a residual
reflection of $0.07$ remains, the $Z_w$--$Z_{\mathrm{eq}}$ gap.  The gap
is closed by measuring $Z$ on the trials themselves: the two seed
evaluations at $a = 0$ and $a = a_{\mathrm{fan}}$ give relaxed states
that differ by one injected frozen wave, and
\begin{equation}
  \ceq^2 = \frac{\tilde P(a_1) - \tilde P(a_0)}{\tilde\rho(a_1) - \tilde\rho(a_0)},
  \qquad Z = \rho_N \ceq ,
  \label{eq:zeq}
\end{equation}
is the equilibrium slope the reaction chain itself imposes on the cell,
whatever the relaxation regime --- no EOS derivative along the
coexistence curve, no neighbour, no model of which legs are fast.  A
degenerate pair keeps the Wood value and is counted.

\subsection{Face types}

The face type is a deck choice per face, never a runtime predicate on
the state.  \code{nr} is the condition above.  \code{vent} is an exit
into a plenum at $\Pinf$, and treats its two flow directions
differently, as every characteristic condition distinguishes them by
the entering waves.  A cell is classed inflow or outflow by the sign of
$\tilde\uout$ on the $a = 0$ trial, the cell's own relaxed evolution
with nothing entering.  An outflow cell takes as its entering wave the
substepped fan of \S3 --- the boundary cell's phase states marched in
pressure toward $\Pinf$ along their isentropes with HEM re-closure
inside the dome, terminated at $\Pinf$ or at the critical-flow point ---
with the wave's amplitude fixed by the fan, no secant, the choke being
the fan's own termination; the trial evaluation is still run on that
wave, so the probe reports the relaxed face state after the fan's wave
has entered.  An inflow cell solves the Hedstrom relation
(\ref{eq:value}), which is the linearised stagnation condition for
fluid entering from a plenum at rest ($P < \Pinf$ by the dynamic head
for inward $u$).  Two earlier forms of the vent were built and refuted
on the way to this one: a static-pressure Dirichlet $\tilde P = \Pinf$
solved by the secant with a choke re-solve $\tilde\uout = \ceq$ where
the ambient solution was relaxed-supersonic, which over-drove the
reversed-flow part of the outlet to supersonic inflow (\S\ref{sec:demo3})
and could not represent the flashing expansion (\S\ref{sec:meas}); and
the same with the plenum inflow relation, which fixed the inflow but
left the outflow cells with a one-jump amplitude that had no root on
the liquid branch on the flashing vent.  The vent type exists because
the production outlet needed it.

\subsection{The scalar solve, refusal, cost}

$g(a)$ is one scalar in one unknown, nearly linear (slope
$\nu\,(\ceq^2/\cf^2 + \ceq/\cf)$ with $\nu$ the swept fraction), solved
by a secant from the seeds $\{0, a_{\mathrm{fan}}\}$ with a fixed number
of updates and no tolerance; the residual is reported.  Any evaluation
that refuses --- a star state with a negative partial mass, a trial
state whose host phase has no EOS root (tested on the EOS's
non-aborting surface with the same dispatch the flux uses), a non-finite
result --- drops the cell to the fan's own wave and is counted.  Three to
four reaction evaluations per boundary cell per step.  There is no Mach
dispatch anywhere: the frozen fan decides whether a wave enters, the
condition decides its strength.

\subsection{Transverse terms and corners}

In 3D-NSCBC the boundary relation is
$\Lin = K(P - \Pinf) + (1-\beta)\mathcal{T}$.  Here the transverse
fluctuations of the boundary cell are part of the known increment
(\ref{eq:split}), so the entering wave compensates them in full ---
$\beta = 0$ --- without anything transverse being computed; $\beta > 0$
would need the transverse change of $I$ separately.  A corner cell is
processed once per face in direction order.

% ---------------------------------------------------------------------
\section{Contrast with Poinsot--Lele}
% ---------------------------------------------------------------------

\begin{description}
\item[What is modelled.]  P\&L model a wave amplitude in a differential
relation and integrate a boundary ODE.  We model the strength of a
discrete wave the interior scheme already propagates; the outgoing waves
are not modelled at all.
\item[What the condition is posed on.]  P\&L impose it on the frozen
eigenstructure.  We impose it on the state after the source has acted,
with the source's own impedance.  Sutherland--Kennedy with relaxation as
the source, without the relaxed eigenstructure on paper.
\item[When it is applied.]  P\&L at the level of the PDE; we at the
level of one time step with the complete increment of the boundary cell
in hand.  Transverse effects enter through that increment.
\item[What the far field means.]  P\&L's non-reflecting limit says the
incoming invariant does not change; the far field enters only through
$K$.  Hedstrom, and the ghost construction we shipped, set the incoming
invariant's \emph{value} to its far-field value every step.  On the
production outlet the difference was decisive (\S\ref{sec:demo3}).
\item[What is deliberately absent.]  A Mach dispatch on any relaxed
speed; any use of the ghost's invariants; any tolerance; any division by
a wave speed --- the near-sonic blow-up of ghost-cell CBC is absent by
construction, its analogue being a conditioning loss of the secant as
the entering wave's swept fraction vanishes.
\end{description}

% ---------------------------------------------------------------------
\section{What has been measured}
\label{sec:meas}
% ---------------------------------------------------------------------

$r$ is reflected over incident pressure amplitude; 256 cells, a 1~m/s
velocity pulse, both faces NSCBC.

\begin{center}
\footnotesize
\setlength{\tabcolsep}{4pt}
\begin{tabular}{lrrrr}
\toprule
case & ghost & relaxed, Wood $Z$ & relaxed, trial $Z$ & wave speed \\
\midrule
N1 vapour 20 bar / 280 K & $+0.0005$ & $-0.00004$ & $-0.00002$ & 234 m/s \\
N1b mixture $\alpha_1 = 0.05$, full closure & $+0.241$ & $+0.067$ & $\mathbf{-0.011}$ & 139 m/s \\
N1b, pressure constraint only & $+0.144$ & $+0.005$ & $+0.005$ & 174 m/s (Wood) \\
N2 2-D pulse, residual/$A$, 4/5/6 ms & .0083/.0067/.0071 & .0091/.0072/.0056 & same & --- \\
\bottomrule
\end{tabular}
\normalsize
\end{center}

The value form on N1 gives the ghost's $+0.0005$ (in vapour the two are
the same condition) and on N1b $+0.0052$.  N4, mean-pressure recovery
from $1.1\Pinf$ through both faces, in units of the initial excess per
millisecond: ghost $1, .763, .528, .293, .060, .008, \ldots$; relaxed
\code{nr} (value) $1, .761, .526, .291, .058, .008, \ldots$; \code{vent}
$1, .54, .08, -.37, -.82, -.73, -.51, -.29, -.07, .007$ (an open end that
rings once and is then absorbed on the inflow phase).  Every run was
bit-identical to the shipped code with the construction off, and no
trial ever had to fall back on the 1-D set.

The flashing vent (the 100~bar supercritical liquid tube of \S3
venting to 1~bar, vented mass at 80~$\mu$s against the boundary-free
reference): the shipped choked fan $+4.4\,\%$; relaxed \code{nr}
$-38\,\%$ (the frozen linear number); the secant \code{vent} $+14\,\%$,
with the secant accepting four early steps at residuals of half the
pressure and refusing every step after; the \code{vent} with the fan as
its outflow wave model reproduces the shipped $+4.4\,\%$ to every digit
with no refusal.  The one-wave form does not do this case; the fan does,
inside the relaxed construction.

\paragraph{Order at the face.}  N1 at 256/512/1024 cells against a
boundary-free reference at the same $\Delta x$ and $\Delta t$ (the same
grid extended so that no face reflection can return in the window),
scoring the face cell's pressure error integrated over the pulse's
passage: the shipped ghost falls by $4.2$ and $3.4$ per refinement, the
relaxed boundary over the same ghost by $4.2$ and $3.4$ --- equal to
four digits --- so the entering wave's dropped correction costs nothing
measurable and the face is second order.  The reflected residual left
in the domain, $0.23$~Pa on a $5$~kPa pulse for every boundary at every
resolution, does not converge: it is the amplitude nonlinearity of the
linearised Hedstrom invariant (a 50~m/s pulse gives $40\times$ the
reflection of a 1~m/s one), not discretisation.

\paragraph{What the ghost still supplies.}  With the entering wave
modelled and the vent's wave the fan's, the ghost state at an NSCBC
face is read only for the outgoing waves at the face and, in reversed
flow, the contact.  Replacing the characteristic ghost by a plain copy
of the boundary cell leaves N1, N4 and the 2-D pulse unchanged, moves
N1b's reflection from $+0.005$ to $-0.006$ (the same magnitude, and
within 10\,\% of the ghost on the boundary-free residual), costs
$1.1\,\%$ of the vented mass at the flashing vent, and drops the face to
first order on outgoing waves (ratios $1.8$, $1.9$; $5$--$22\times$ the
face-cell error): what a copy removes is the outgoing wave's
second-order correction, not anything characteristic.  Replacing it
instead by a minmod continuation of the phase states past the face
(the mixture assembled from the phases so that $\rho = m_1 + m_2$ and
$\rho E = \mathcal{E}_1 + \mathcal{E}_2$ hold exactly; a continuation
with no EOS root refused to the copy and counted) recovers the ghost's
order and values to four digits on N1 and N4, $-0.28\,\%$ on the
flashing vent, and the smallest residual of the three on N1b.  If that
holds on the production outlet, the characteristic construction ---
invariants, extrapolation, Mach dispatch, strain cap --- has nothing
left to supply and is retired.

\subsection{The production outlet}
\label{sec:demo3}

A saturated two-phase jet from a pipe rupture reaches the outflow plane
20 diameters downstream as an overexpanded supersonic core
($u/c \approx 1.5$, $13$~bar in gas at $20$~bar) with 60--70\,\% of the
face in reversed flow.  The shipped ghost relieves the arrival transient
within a millisecond and lets the core leave supersonically.  Each
version of the relaxed construction was restarted from the same
checkpoint before arrival and run through it.

\begin{description}
\item[Rate form, \code{nr}.]  Robust --- no refusal on any face in 3000
steps, the validator clean --- and wrong: the last cells hold 40--47~bar
against a supersonic core behind them.  The probe shows the ghost
sending an expansion of $-0.5$ to $-2$~MPa per step at arrival and the
relaxed condition choosing $+40$~kPa.  $\Lin = 0$ carries no far field:
once the arrival raised $P$ and the flow reversed, $P - Z\uout$ stayed
where it was.
\item[Value form, \code{nr}.]  Worse, 46--58~bar and the shock ten cells
inside by 24~ms, for a reason that took the numbers to see: at the
arrival cell $P = 38$~bar, $u = 85$~m/s, $\rho c \approx 3.5\times10^4$,
so $P - \rho c u \approx 8$~bar, \emph{below} $\Pinf$, and the Hedstrom
invariant condition asks for a compression.  The ghost's linear branch
would too; the ghost relieves because the jump exceeds its strain cap
and it switches to the choked fan, which ignores the invariants.  The
shipped face is a characteristic outflow at small jumps and a venting
outlet at large ones, with a threshold deciding.  An acoustic
linearisation about a quiescent far field accepts $P \ne \Pinf$
compensated by velocity; an exit into a plenum does not.
\item[\code{vent}, Dirichlet on all subsonic cells.]  The arrival is
relieved (harder than the ghost: face at 6.5~bar, core through at
$u/c = 1.95$) and the picture at $+2$~ms matches the ghost's.  Later it
goes the other way: the reversed-flow part of the face is driven to
supersonic \emph{inflow} ($u/c$ to $-2.2$), the mass pushed in
recompresses the core and the shock walks inside.  A static-pressure
Dirichlet is the exit condition for outflow into a plenum; for inflow
from a plenum at rest the condition is stagnation.
\item[\code{vent} as a plenum (inflow on the Hedstrom relation).]  The
arrival as before; the reversed-flow pathology gone (inflow $u/c$ no
lower than $-0.16$), the inflow cells sitting below $\Pinf$ as the
stagnation relation has them, the time step tracking the ghost's.  The
core takes a detour the ghost does not --- subsonic at 22~bar at
$+4$~ms where the ghost is supersonic at 14, supersonic again at
$+8$~ms --- a recompression excursion consistent with the open-end ring
of the N4 test acting on the subsonic surroundings; whether that or the
ghost's earlier settling is the better outlet is not decidable against
the ghost, since both are boundary constructions.  Two to three cells
per step still refused to the fan's wave in the late phase.
\item[\code{vent} with the fan as the outflow wave model.]  Identical to
the previous run to the second decimal at every station; the refusals
fall from 7979 in 3000 steps to 1, because the fan reaches the same
relaxed state along an admissible path with a wave one tenth the
secant's amplitude.  This is the configuration that ships: \code{nr}
on the lateral faces, \code{vent} on the outlet, the value form, the
trial impedance.
\item[Ghost fill: copy.]  With the characteristic ghost replaced by a
copy of the boundary cell, the outlet is again identical to the
previous run to the second decimal through the reversed-flow phase,
with no refusal on any face: the plenum inflow relation on the cell,
not the ghost's contact wave, is what sets the inflow.  The one change
is diagnostic --- with no wave sent by the ghost the two seed trials of
(\ref{eq:zeq}) coincide on a third of the lateral-face cells and the
impedance falls back to Wood there, harmlessly (the cells are at rest
near $\Pinf$).
\item[Ghost fill: continuation.]  A first form that continued every
conserved slot separately was refuted by the validator on the first
step (a hundred boundary cells with $\rho \ne m_1 + m_2$ at $10^{-9}$)
and, 2300 steps later, by a time-step collapse and an abort; the
assembled form of \S\ref{sec:meas} is running at the time of writing.
\end{description}

% ---------------------------------------------------------------------
\section{Questions for discussion}
% ---------------------------------------------------------------------

\begin{enumerate}
\item \textbf{Value or rate.}  Hedstrom's value condition on the relaxed
invariant against Poinsot--Lele's rate condition with a $K$.  On plane
waves both are transparent; on the impinging outlet the value form
anchors strongly and still could not vent.  Is there a principled way to
choose the anchoring strength for a face that sees both small-amplitude
acoustics and an order-one transient, other than a threshold on the
jump?
\item \textbf{The source along the wave.}  For a flashing expansion the
relaxation acts along the fan, not after the wave has entered the cell,
and a frozen star jump cannot represent it.  Is there an analogue in
reacting flow --- a case where the reaction must be integrated along the
characteristic rather than in the boundary relation --- and how was it
treated?  Our answer is to keep the substepped fan as the outflow
wave model of the vent type and use the trial evaluation only to read
the relaxed result; it reproduces the flashing-vent benchmark exactly
and the production outlet to the second decimal.
\item \textbf{Reversed flow.}  On the part of the face where the flow
enters, the acoustic wave is the modelled one and the contact wave
carries whatever the ghost holds.  What should enter: the cell's
own fluid, or the far-field composition and temperature?  On our outlet
the answer made no measurable difference --- a ghost that sends no
contact at all reproduces the outlet, because the plenum inflow
relation on the cell sets the inflow --- but the far field there has
the composition of the jet.  In reacting flow this is the burnt-gas
backflow question; the answer probably carries over.
\item \textbf{The impedance from the trials.}  Equation (\ref{eq:zeq})
measures the equilibrium slope the source imposes over one step and
adapts to whatever the relaxation regime is.  Is there a cleaner
statement of what invariant a relaxing system conserves across a
subsonic outflow, so that $Z$ need not be measured?
\item \textbf{Transverse terms.}  Compensating the boundary cell's
transverse fluctuations in full ($\beta = 0$) through the known
increment left the 2-D pulse's corner residual unchanged at about
$6\,\%$.  Is the Yoo--Im $\beta = M$ compromise still the recommendation
when the transverse terms are known exactly rather than estimated?
\item \textbf{Order and conditioning.}  The modelled wave enters first
order (no limited correction for a wave with no upwind neighbour), yet
the measured order at the face is the ghost's, second, because the
entering wave on an outflow is small; what does cost an order is a ghost
that carries no outgoing wave (a copy).  Masset et al.'s bound --- the
minimum over the extrapolation, the outgoing-wave estimate and the
scheme --- reads the same way.  Is there a case in reacting flow where
the entering wave is large enough for its order to matter, and what was
done there?
\item \textbf{The relaxed-supersonic reading.}  The frozen fan decides
whether a wave enters; when the flow is relaxed-supersonic but
frozen-subsonic the condition is still imposed.  A dispatch on the
relaxed Mach number was refuted on this outlet (it held the jet head
against the plane).  Is there a better statement for that regime?
\end{enumerate}

% ---------------------------------------------------------------------
\section*{References}
% ---------------------------------------------------------------------
\begin{itemize}
\item T.~J.~Poinsot, S.~K.~Lele, Boundary conditions for direct simulations of compressible viscous flows, \emph{J.\ Comput.\ Phys.} 101 (1992).
\item K.~W.~Thompson, Time dependent boundary conditions for hyperbolic systems, \emph{J.\ Comput.\ Phys.} 68 (1987).
\item G.~W.~Hedstrom, Nonreflecting boundary conditions for nonlinear hyperbolic systems, \emph{J.\ Comput.\ Phys.} 30 (1979).
\item D.~H.~Rudy, J.~C.~Strikwerda, A nonreflecting outflow boundary condition for subsonic Navier--Stokes calculations, \emph{J.\ Comput.\ Phys.} 36 (1980).
\item J.~C.~Sutherland, C.~A.~Kennedy, Improved boundary conditions for viscous, reacting, compressible flows, \emph{J.\ Comput.\ Phys.} 191 (2003).
\item C.~S.~Yoo, H.~G.~Im, Characteristic boundary conditions for simulations of compressible reacting flows with multi-dimensional, viscous and reaction effects, \emph{Combust.\ Theory Model.} 11 (2007).
\item G.~Lodato, P.~Domingo, L.~Vervisch, Three-dimensional boundary conditions for direct and large-eddy simulation of compressible viscous flows, \emph{J.\ Comput.\ Phys.} 227 (2008).
\item E.~Motheau, A.~Almgren, J.~B.~Bell, Navier--Stokes characteristic boundary conditions using ghost cells, \emph{AIAA J.} 55 (2017).
\item P.-A.~Masset, L.~Lecointre, C.~Tenaud, M.~Massot, S.~Koudriakov, E.~Studer, Characteristic boundary conditions using ghost cells (GC-CBC) for high-order finite volume simulations of multicomponent Euler gas equations handling near-zero/sonic velocities, preprint.
\item M.~Pelanti, K.-M.~Shyue, A mixture-energy-consistent six-equation two-phase numerical model for fluids with interfaces, cavitation and evaporation waves, \emph{J.\ Comput.\ Phys.} 259 (2014).
\item R.~J.~LeVeque, \emph{Finite Volume Methods for Hyperbolic Problems}, Cambridge (2002).
\end{itemize}

\end{document}
